\documentclass[11pt,a4paper]{article}

\usepackage[utf8]{inputenc}
\usepackage[T1]{fontenc}
\usepackage[english]{babel}
\usepackage{xcolor}
\usepackage{amsmath,amssymb,amsthm}
\usepackage{mathtools}
\usepackage{geometry}
\geometry{margin=2.4cm}
\usepackage{booktabs}
\usepackage{tabularx}
\usepackage{enumitem}
\usepackage{tikz-cd}
\usepackage[most]{tcolorbox}
\usepackage{hyperref}
\hypersetup{
  colorlinks=true,
  linkcolor=blue!60!black,
  citecolor=blue!60!black,
  urlcolor=blue!60!black
}

% theorem environments
\newtheorem{theorem}{Theorem}[section]
\newtheorem{proposition}[theorem]{Proposition}
\newtheorem{corollary}[theorem]{Corollary}
\newtheorem{lemma}[theorem]{Lemma}
\newtheorem{conjecture}[theorem]{Conjecture}
\newtheorem{ansatz}[theorem]{Phenomenological Ansatz}
\theoremstyle{definition}
\newtheorem{definition}[theorem]{Definition}
\newtheorem{axiom}{Axiom}
\theoremstyle{remark}
\newtheorem{remark}[theorem]{Remark}

% shortcuts
\newcommand{\Sn}[1]{S^{#1}}
\newcommand{\Szero}{S^{0}}
\newcommand{\Sone}{S^{1}}
\newcommand{\Stwo}{S^{2}}
\newcommand{\Sthree}{S^{3}}
\newcommand{\Sfour}{S^{4}}
\newcommand{\Sfive}{S^{5}}
\newcommand{\Sseven}{S^{7}}
\newcommand{\CP}{\mathbb{C}P}
\newcommand{\HP}{\mathbb{H}P}
\newcommand{\R}{\mathbb{R}}
\newcommand{\C}{\mathbb{C}}
\newcommand{\Hbb}{\mathbb{H}}
\newcommand{\Z}{\mathbb{Z}}
\newcommand{\SU}{\mathrm{SU}}
\newcommand{\SO}{\mathrm{SO}}
\newcommand{\U}{\mathrm{U}}
\newcommand{\Spin}{\mathrm{Spin}}
\newcommand{\Vol}{\mathrm{Vol}}
\newcommand{\Id}{\mathrm{id}}
\newcommand{\golden}{\varphi}
\newcommand{\dd}{\mathrm{d}}
\newcommand{\coker}{\mathrm{coker}}
\newcommand{\rank}{\mathrm{rank}}

\title{\bfseries Suspension of Spheres and the Minimal Role of \texorpdfstring{$\Sthree$}{S3}\\
in Spin, Hopf Geometry, and Irreversible Hierarchy}

\author{%
\texorpdfstring{$^{1}$}{1}Andrei Preece\thanks{\texttt{andre.preece@web.de}; ORCID: 0009-0004-5978-795X}
\quad and \quad
\texorpdfstring{$^{1}$}{1}Boris Batenin\thanks{\texttt{batenin@web.de}; ORCID: 0009-0000-7027-8583}\\[0.35em]
\small \texorpdfstring{$^{1}$}{1}Independent Researchers
}

\date{Berlin, 2026}

\begin{document}
\maketitle

\begin{abstract}
We reconstruct the suspension-chain argument
\[
\Szero \longrightarrow \Sone \longrightarrow \Stwo \longrightarrow \Sthree \longrightarrow \Sfour
\]
as a mathematical-core paper in topology, spin geometry, Hopf geometry, and scale hierarchy.
The claim is deliberately narrower than a complete cosmology.  The result is not that topology
alone derives matter content, amplitudes, or a final theory of physics.  The result is that
$\Sthree$ occupies a distinguished minimal position in the suspension ladder: it is the first
sphere which is simultaneously compact, simply connected, parallelizable, a non-Abelian Lie
group, the total space of the nontrivial Hopf fibration
\[
\Sone \hookrightarrow \Sthree \twoheadrightarrow \Stwo,
\]
and a natural carrier of spinorial and phase-direction geometry.  The paper proves this carrier
claim, separates it from physical over-interpretation, and formulates the remaining hierarchy
claims as geometric consequences, phenomenological ansätze, or conjectures.

The four-sphere is treated not as a replacement for physical three-dimensional spatial reality,
but as a compensation endpoint in the suspension chain: $\Sfour\simeq\HP^{1}$ carries the
quaternionic projective completion of the preceding $\Sthree$ structure, while the local
spin geometry remains governed by $\SU(2)$ and its two chiral copies inside
$\SO(4)\simeq(\SU(2)_{L}\times\SU(2)_{R})/\Z_{2}$.  The golden-ratio hierarchy is retained
as a minimal additive self-similar scaling law; its entropy monotonicity is proved under an
area/holographic counting assumption, while recurrence breaking is explicitly left as a
conjecture requiring a specified Hamiltonian measure.  Observational consequences such as
low-multipole suppression and log-periodic power-spectrum modulation are presented only as
phenomenological tests, not as evidence for topology.  A new section develops the lowest level
of the chain: $\Szero$ as minimal interdependent distinction, its non-contractibility within
structure-preserving operations, the transport of the distinction class onto the fundamental
class of $\Sthree$ by the suspension isomorphism, and the Planck scale as the lower operational
realization of distinguishability; all interpretive readings are labelled as such.

\medskip
\noindent\textbf{Keywords:} suspension of spheres; $\Sthree$ topology; $\SU(2)$; Hopf fibration;
Chern class; spin structure; parallelizability; golden-ratio hierarchy; compact cosmology;
irreversible hierarchy.
\end{abstract}

\tableofcontents

\section{Introduction}

The purpose of this paper is to isolate the mathematical core of an idea that is easily
overstated if presented as cosmology before it is presented as topology.  The core question is not
``why does the universe exist?'' and not ``can topology alone explain all physics?''  The precise
question is the following.

\begin{quote}
What is the minimal compact sphere in the suspension hierarchy that can simultaneously carry
phase, direction, spin, non-Abelian internal structure, nowhere-vanishing flow, and a nontrivial
phase--direction bundle?
\end{quote}

The answer developed here is $\Sthree$.  The suspension ladder
\begin{equation}
\Szero \xrightarrow{\Sigma} \Sone \xrightarrow{\Sigma} \Stwo
\xrightarrow{\Sigma} \Sthree \xrightarrow{\Sigma} \Sfour
\label{eq:suspension-ladder}
\end{equation}
is not used as decorative ontology.  It is used as a bookkeeping device for the following
structural transition:
\[
\begin{aligned}
\text{binary distinction}
&\longrightarrow \text{phase}
\longrightarrow \text{direction}\\
&\longrightarrow \text{spin/volume/Hopf total space}
\longrightarrow \text{quaternionic compensation}.
\end{aligned}
\]
The special role of $\Sthree$ is not that one may pronounce it philosophically important.  Its
special role is that the relevant structures meet there for the first time.  We use
\emph{carrier} in this restricted sense: a carrier is a manifold on which the named structures
exist naturally and compatibly.  It is not, by that fact alone, a derivation of dynamics,
matter content, amplitudes, or observational data.

\subsection{Claim-status convention}

The paper uses five status labels.

\begin{center}
\small
\begin{tabularx}{\textwidth}{@{}lX@{}}
\toprule
\textbf{Label} & \textbf{Meaning} \\
\midrule
Theorem & Proved under stated mathematical assumptions.\\
Geometric consequence & Standard consequence of topology, geometry, or representation theory.\\
Ansatz & A proposed scaling or phenomenological form, not derived from dynamics.\\
Conjecture & A statement requiring additional Hamiltonian, measure, or dynamical data.\\
Interpretation & Philosophy-of-physics reading; not a proof of physics.\\
\bottomrule
\end{tabularx}
\end{center}

This convention is essential.  The paper preserves the philosophical force of the suspension
picture, but it does not promote philosophical interpretation to theorem.

\subsection{What this paper does not claim}

\begin{tcolorbox}[colback=gray!4,colframe=gray!50,title=Boundary of the paper]
This paper does not prove that the observed universe is globally $\Sthree$; does not derive the
Standard Model; does not derive dark matter or dark energy abundances; does not derive the
numerical value of any fundamental constant; does not derive the spin-statistics theorem from
topology alone; and does not treat the golden ratio as a fundamental physical coupling.
\end{tcolorbox}

\begin{tcolorbox}[colback=blue!3,colframe=blue!45!black,title=Structural roadmap]
The argument has the following invariant spine:
\[
\Szero \xrightarrow{\Sigma} \Sone \xrightarrow{\Sigma} \Stwo
\xrightarrow{\Sigma} \Sthree \xrightarrow{\Sigma} \Sfour ,
\]
with
\[
\begin{gathered}
\text{distinction}\to\text{phase}\to\text{direction}\\
\to\text{Hopf total space/spin carrier}
\to\text{quaternionic compensation}.
\end{gathered}
\]
The proof does not proceed from analogy.  It proceeds by checking which level first carries
the joint package
\[
(\pi_{1}=0,\;\SU(2),\;c_{1}\neq0,\;w_{2}=0,\;T\Sthree\simeq\Sthree\times\R^{3}).
\]
\end{tcolorbox}

The positive claim is narrower and stronger:
\[
\begin{aligned}
\Sthree\quad &\text{is the minimal compact simply connected spherical carrier}\\
&\text{of the joint }(\SU(2),\text{ Hopf},\text{ spin},\text{ parallelizable flow})
\text{ structure}.
\end{aligned}
\]

\section{Suspension Ladder and Sphere Hierarchy}

\begin{definition}[Suspension]
For a topological space $X$, its suspension is
\begin{equation}
\Sigma X := (X\times[0,1])/\sim,
\end{equation}
where $(x,0)\sim N$ and $(x,1)\sim S$ for all $x\in X$.  The two equivalence classes
$N$ and $S$ are the suspension poles.
\end{definition}

For spheres one has the standard homeomorphism
\begin{equation}
\Sigma \Sn{n}\simeq \Sn{n+1}.
\label{eq:sigma-sphere}
\end{equation}
Starting from $\Szero=\{-1,+1\}$, this gives the chain \eqref{eq:suspension-ladder}.

\begin{center}
\small
\begin{tabularx}{\textwidth}{@{}c l l X@{}}
\toprule
Level & Sphere & Main structure & Minimal role in the chain\\
\midrule
0 & $\Szero$ & two points & binary distinction\\
1 & $\Sone\simeq\U(1)$ & Abelian Lie group & phase\\
2 & $\Stwo\simeq\CP^{1}$ & complex projective line & direction/Bloch sphere\\
3 & $\Sthree\simeq\SU(2)$ & non-Abelian Lie group, Hopf total space & spin, phase--direction unity, volume\\
4 & $\Sfour\simeq\HP^{1}$ & quaternionic projective line & compensation endpoint\\
\bottomrule
\end{tabularx}
\end{center}


The word ``minimal'' is used only in the order of the chain \eqref{eq:suspension-ladder}.  It
does not assert uniqueness among all manifolds; it asserts first occurrence among the suspended
spheres under discussion.

\begin{proposition}[Interdependence of the middle levels]
The Hopf fibration
\begin{equation}
\Sone \hookrightarrow \Sthree \xrightarrow{p} \Stwo
\label{eq:hopf}
\end{equation}
uses $\Sone$ as fibre, $\Stwo$ as base, and $\Sthree$ as total space.  Thus phase and
direction are not joined by a trivial Cartesian product $\Stwo\times\Sone$; they are joined by
the nontrivial bundle \eqref{eq:hopf}.
\end{proposition}

\begin{proof}
The classical Hopf map $p:\Sthree\to\Stwo$ is a principal $\U(1)$-bundle.  Its associated
complex line bundle has first Chern class $c_{1}=1\in H^{2}(\Stwo;\Z)$, so it is not isomorphic
to the trivial bundle.  Therefore $\Sthree$ is not globally $\Stwo\times\Sone$ as a principal
$\U(1)$-bundle over $\Stwo$.
\end{proof}

The suspension hierarchy is therefore not merely dimensional.  It records the point at which
phase and direction acquire a nontrivial global total space.

\subsection{Geometric intuition: what suspension changes}

Suspension is a topological operation, not merely an index shift.  The quotient
\[
\Sigma X=(X\times[-1,1])/(X\times\{-1\}\sim N,\;X\times\{1\}\sim S)
\]
keeps the local product structure away from the two poles, while concentrating the link of
$X$ into the neighbourhoods of $N$ and $S$.  If $X$ is already a smooth sphere, this produces
the ordinary smooth sphere one dimension higher.  If $X$ is only a homology sphere, the first
suspension may have pole neighbourhoods whose links are homology spheres rather than standard
spheres; the resulting object is then topologically meaningful but not automatically a smooth or
PL manifold at the poles.

The double-suspension theorem clarifies why this distinction matters.  For a homology sphere
$H^{n}$, the double suspension $\Sigma^{2}H^{n}$ is homeomorphic to $\Sn{n+2}$.  Thus a second
suspension can remove a topological local obstruction at the level of homeomorphism type.
This is not the same as preserving a PL structure, nor is it a statement about exotic smooth
structures.  Suspension may improve the topological manifold status while losing PL rigidity.
The lesson for the present chain is simple: the ladder records topological carrier structure first;
smooth and PL claims must be stated separately.

\begin{definition}[Exotic sphere]
An exotic $n$-sphere is a smooth manifold homeomorphic to the standard sphere $\Sn{n}$ but
not diffeomorphic to it.  In this paper the word ``exotic'' is used only in this smooth sense.
A topological sphere, a PL sphere, and a smooth sphere are therefore not interchangeable
objects unless the relevant equivalence category has been specified.
\end{definition}

\begin{remark}[Category discipline]
The suspension chain in this paper is a topological carrier chain.  Whenever smooth spin
geometry, Dirac spectra, or Lie-group structure are used, the smooth structure is the standard
smooth structure on the corresponding round sphere unless explicitly stated otherwise.
\end{remark}

\section{Minimality of \texorpdfstring{$\Sthree$}{S3}}

\subsection{Lie-group spheres}

Among spheres, the associative Lie group cases are
\begin{equation}
\Sone\simeq \U(1),\qquad \Sthree\simeq\SU(2).
\end{equation}
The seven-sphere is related to unit octonions, but octonionic multiplication is non-associative;
therefore $\Sseven$ is not a Lie group in the strict associative sense.

The explicit diffeomorphism $\Sthree\to\SU(2)$ is
\begin{equation}
(z_{1},z_{2})\longmapsto
\begin{pmatrix}
z_{1} & -\overline{z}_{2}\\
z_{2} & \overline{z}_{1}
\end{pmatrix},
\qquad |z_{1}|^{2}+|z_{2}|^{2}=1.
\label{eq:s3-su2}
\end{equation}
This is the first sphere in the suspension ladder that is both simply connected and a non-Abelian
Lie group.

\subsection{Hopf nontriviality and Chern class}

\begin{lemma}[Hopf bundle]
The map
\begin{equation}
p:\Sthree\to\Stwo,\qquad
p(z_{1},z_{2})=
\left(
2z_{1}\overline{z}_{2},
|z_{1}|^{2}-|z_{2}|^{2}
\right)
\in \C\times\R\simeq\R^{3}
\label{eq:hopf-map}
\end{equation}
defines a principal $\U(1)$-bundle
\[
\Sone\hookrightarrow\Sthree\xrightarrow{p}\Stwo
\]
with first Chern class
\begin{equation}
c_{1}=1\in H^{2}(\Stwo;\Z)\simeq\Z.
\end{equation}
\end{lemma}

\begin{proof}
The fibre action is $(z_{1},z_{2})\mapsto(e^{i\alpha}z_{1},e^{i\alpha}z_{2})$, which leaves
\eqref{eq:hopf-map} invariant.  The homotopy class of $p$ generates
$\pi_{3}(\Stwo)\simeq\Z$, and the associated principal $\U(1)$-bundle over $\Stwo$ is classified
by its first Chern class.  For the Hopf bundle this class is the generator of
$H^{2}(\Stwo;\Z)$.
\end{proof}

\subsection{Spin structure and Dirac spectrum}

Every orientable three-manifold is spin; in particular
\begin{equation}
w_{2}(T\Sthree)=0.
\end{equation}
For the round $\Sthree_{R}$, the Dirac operator has discrete spectrum
\begin{equation}
\mathrm{spec}(D_{\Sthree_{R}})
=
\left\{
\pm\frac{1}{R}\left(n+\frac{3}{2}\right): n=0,1,2,\ldots
\right\},
\label{eq:dirac-spectrum}
\end{equation}
with finite multiplicities.  This discrete spectrum follows from compactness and ellipticity; it is
not a lattice cutoff.

\begin{remark}[Spin structure versus spin-statistics]
The identification $\Sthree\simeq\SU(2)$ gives the natural double cover
\[
\SU(2)\to\SO(3),
\]
and therefore the geometric carrier for spin-$1/2$ representations.  The full spin-statistics
theorem additionally requires the axioms of relativistic quantum field theory, including locality
and positivity conditions.  Topology supplies the carrier; it is not by itself the complete theorem.
\end{remark}

\subsection{Parallelizability and nowhere-vanishing flows}

Because $\Sthree\simeq\SU(2)$ is a Lie group, its tangent bundle is trivial:
\begin{equation}
T\Sthree\simeq \Sthree\times\R^{3}.
\end{equation}
Equivalently, $\Sthree$ is parallelizable.  In complex coordinates $\Sthree\subset\C^{2}$, an
explicit nowhere-vanishing tangent vector field is
\begin{equation}
V(z_{1},z_{2})=(iz_{1},iz_{2}).
\label{eq:hopf-vector-field}
\end{equation}

\begin{theorem}[Minimality of $\Sthree$ as geometric carrier]
Among the spheres in the suspension chain, $\Sthree$ is the lowest-dimensional sphere that
simultaneously satisfies:
\begin{enumerate}[label=(C\arabic*)]
\item compactness and simple connectivity;
\item associative non-Abelian Lie-group structure;
\item nontrivial principal $\U(1)$ Hopf fibration over $\Stwo$;
\item vanishing $w_{2}$ and a natural spin carrier;
\item parallelizability and nowhere-vanishing smooth vector fields.
\end{enumerate}
\end{theorem}

\begin{proof}
$\Szero$ is disconnected and has no manifold Lie-group structure relevant to spin geometry.
$\Sone$ is a compact Lie group, but it is Abelian and not simply connected.  $\Stwo$ is simply
connected but admits no Lie-group structure and is not parallelizable; indeed $\chi(\Stwo)=2$,
so every smooth tangent vector field has a zero.  $\Sthree$ is compact, simply connected,
parallelizable, and diffeomorphic to the non-Abelian Lie group $\SU(2)$; it is also the total
space of the Hopf fibration with $c_{1}\neq0$.  $\Sfour$ is higher-dimensional and is not a Lie
group.  Thus $\Sthree$ is the first sphere in the chain satisfying all five conditions.
\end{proof}

\subsection{Contrast with \texorpdfstring{$\Stwo$}{S2}}

The two-sphere is geometrically indispensable as the base of the Hopf fibration, but it is not
sufficient as a carrier of the full structure.

\begin{proposition}[Three failures of $\Stwo$ as total carrier]
The sphere $\Stwo$ fails to be the total carrier of the spin--Hopf--flow structure because:
\begin{enumerate}[label=(\roman*)]
\item $\chi(\Stwo)=2$, so every smooth tangent vector field has a zero;
\item $\Stwo$ admits no Lie-group structure;
\item $\Stwo$ is the base, not the total space, of the nontrivial Hopf bundle.
\end{enumerate}
\end{proposition}

\begin{proof}
The first statement is the hairy-ball theorem, or equivalently the Poincaré--Hopf theorem.
The second follows from the classification of Lie-group spheres.  The third follows from the
bundle sequence $\Sone\hookrightarrow\Sthree\twoheadrightarrow\Stwo$.
\end{proof}

\section{Dynamical Consequences}

\subsection{Intrinsic versus externally imposed motion}

A Lie group carries canonical left- and right-invariant vector fields.  For
$\Sthree\simeq\SU(2)$, these fields are internal to the manifold rather than imposed by an
external embedding.  This is the precise mathematical sense in which $\Sthree$ supports
intrinsic dynamics.

Let $\{e_{1},e_{2},e_{3}\}$ be a basis of $\mathfrak{su}(2)$ with
\begin{equation}
[e_{i},e_{j}]=2\epsilon_{ijk}e_{k}.
\end{equation}
Left translation gives a global frame on $\SU(2)$:
\begin{equation}
E_{i}(g)=(L_{g})_{*}e_{i},\qquad g\in\SU(2).
\end{equation}
Thus the local algebra is non-Abelian at every point, but the frame is globally defined.

\subsection{The role of \texorpdfstring{$\Sfour$}{S4}}

The four-sphere appears in the suspension chain as
\[
\Sigma\Sthree\simeq\Sfour.
\]
It also has the projective representation
\begin{equation}
\HP^{1}\simeq\Sfour.
\end{equation}
This makes $\Sfour$ a natural compensation endpoint of the chain in a restricted sense: the quaternionic line
contains the preceding $\SU(2)$ phase structure as fibre data in the quaternionic Hopf map
\[
\Sthree\hookrightarrow\Sseven\twoheadrightarrow\Sfour.
\]

The relevant symmetry relation is
\begin{equation}
\Spin(4)\simeq \SU(2)_{L}\times\SU(2)_{R},
\qquad
\SO(4)\simeq(\SU(2)_{L}\times\SU(2)_{R})/\Z_{2}.
\label{eq:so4}
\end{equation}

\begin{definition}[Chirality compensation]
A pair $(U_{L},U_{R})\in\SU(2)_{L}\times\SU(2)_{R}$ is called compensated if
\begin{equation}
U_{R}=U_{L}^{-1}.
\end{equation}
In the quotient \eqref{eq:so4}, such pairs encode cancellation of opposite left/right rotations
in the corresponding $\SO(4)$ action.
\end{definition}

\begin{remark}
The statement that $\Sfour$ represents a compensation endpoint is an interpretation of the
suspension and $\Spin(4)$ structure.  It is not a theorem that physical reality becomes $\Sfour$,
nor that $\Sfour$ replaces $\Sthree$ as spatial carrier.
\end{remark}

\subsection{Why further suspension is not used as new physical content}

Mathematically, $\Sigma\Sfour\simeq\Sfive$ exists.  The claim made here is not that $\Sfive$
is impossible.  The claim is that the suspension chain has already passed through the minimal
spin--Hopf--volume carrier at $\Sthree$ and the quaternionic compensation endpoint at $\Sfour$.
A further suspension adds a new topological sphere, but it does not add a new member to the
specific chain of binary distinction, phase, direction, total Hopf volume, and compensation used
in this paper.

\section{\texorpdfstring{$\golden$}{phi}-Scaling and Entropy Hierarchy}

\subsection{Minimal additive self-similarity}

The golden-ratio hierarchy is kept as a structural scaling ansatz, not as a fundamental coupling.

\begin{axiom}[Self-similarity]
There is a scale sequence $R_{n}>0$ such that
\begin{equation}
R_{n+1}=\alpha R_{n},\qquad \alpha>1.
\end{equation}
\end{axiom}

\begin{axiom}[Additive closure]
Adjacent levels close by
\begin{equation}
R_{n+1}=R_{n}+R_{n-1}.
\end{equation}
\end{axiom}

\begin{theorem}[Golden-ratio selection]
The two axioms imply
\begin{equation}
\alpha=\golden:=\frac{1+\sqrt5}{2}.
\end{equation}
\end{theorem}

\begin{proof}
Using $R_{n-1}=R_{n}/\alpha$, the additive closure axiom gives
\[
\alpha R_{n}=R_{n}+\frac{R_{n}}{\alpha}.
\]
Dividing by $R_{n}$ and multiplying by $\alpha$ yields
\[
\alpha^{2}=\alpha+1.
\]
The positive root greater than one is $\golden=(1+\sqrt5)/2$.
\end{proof}

Thus
\begin{equation}
R_{n}=R_{0}\golden^{n}.
\label{eq:phi-ladder}
\end{equation}

\begin{remark}[Status of $\golden$]
Equation \eqref{eq:phi-ladder} is a minimal self-similar additive hierarchy.  It is not, by itself,
a derivation of cosmic amplitudes.  Any claim that the hierarchy is physically selected requires
an additional dynamical selection principle.
\end{remark}

\subsection{No exact multiplicative period}

\begin{lemma}[Aperiodicity]
The hierarchy $R_{n}=R_{0}\golden^{n}$ has no nontrivial exact multiplicative period:
\[
\golden^{k}=1\quad\Longrightarrow\quad k=0.
\]
\end{lemma}

\begin{proof}
Since $\golden>1$, $\golden^{k}>1$ for $k>0$ and $0<\golden^{k}<1$ for $k<0$.  Hence only
$k=0$ gives $\golden^{k}=1$.
\end{proof}

This elementary aperiodicity is weaker than a proof of thermodynamic irreversibility.  It only
states that the scaling hierarchy does not close by an exact multiplicative cycle.

\subsection{Area-counting entropy}

For a compact three-sphere with radius $R_{n}$, the equatorial two-sphere area scale is
\begin{equation}
A_{n}=4\pi R_{n}^{2}=4\pi R_{0}^{2}\golden^{2n}.
\end{equation}
If one uses an area/holographic counting rule
\begin{equation}
S_{n}=\eta\, k_{B}\frac{A_{n}}{\ell_{P}^{2}},
\qquad \eta>0,
\label{eq:entropy-area}
\end{equation}
then the entropy grows monotonically with $n$.

\begin{theorem}[Monotonic hierarchy entropy]
Under the counting rule \eqref{eq:entropy-area},
\begin{equation}
\frac{\dd S_{n}}{\dd n}
=
2(\ln\golden)S_{n}>0.
\end{equation}
\end{theorem}

\begin{proof}
Substitute $A_{n}=4\pi R_{0}^{2}\golden^{2n}$ into \eqref{eq:entropy-area}.  Differentiation with
respect to the continuous interpolation of $n$ gives the stated result.
\end{proof}

\begin{remark}[No volume-scaling shortcut]
The area-counting rule is chosen because entropy bounds in gravity are area-based.  A volume
formula $\exp(V/\ell_{P}^{3})$ is not used as a theorem here.
\end{remark}

\subsection{Recurrence status}

\begin{conjecture}[Effective non-recurrence of an open hierarchy]
If the hierarchy \eqref{eq:phi-ladder} is promoted to an unbounded sequence of dynamically
accessible sectors with a measure whose total phase-space volume diverges, then ordinary
Poincaré recurrence need not apply to the full hierarchy.
\end{conjecture}

\begin{remark}
This is not a theorem until the Hamiltonian, state space, invariant measure, and allowed
transitions between levels are specified.  Each fixed compact level may still satisfy recurrence
hypotheses under an appropriate finite invariant measure.
\end{remark}

\section{Relation to Bell Geometry and Compact Quantum State Space}

The suspension result is useful for quantum geometry because the Hopf fibration is already the
geometry of a single qubit.  Pure one-qubit states form
\[
\CP^{1}\simeq\Stwo,
\]
while normalized spinors lie on $\Sthree\subset\C^{2}$:
\[
|\psi\rangle =
\begin{pmatrix}
z_{1}\\ z_{2}
\end{pmatrix},
\qquad |z_{1}|^{2}+|z_{2}|^{2}=1.
\]
The projection
\[
\Sthree\to\CP^{1}\simeq\Stwo
\]
quotients by the physically irrelevant global $\U(1)$ phase.  The same map is the Hopf map.

\begin{proposition}[Phase-direction non-factorization]
The normalized spinor space $\Sthree$ is not globally the product $\Stwo\times\Sone$ as a
principal $\U(1)$-bundle over $\Stwo$.
\end{proposition}

\begin{proof}
This is exactly the nontriviality of the Hopf bundle, expressed by $c_{1}=1$.
\end{proof}

\begin{remark}[Bell correlations]
Bell-inequality violation is not derived in this paper.  What is established here is the geometric
carrier: the nontrivial Hopf bundle that prevents the normalized spinor from being globally
decomposed into an independent direction variable and an independent phase variable.  Any Bell
analysis must still specify measurement maps, probability rules, and locality assumptions.
\end{remark}

\section{Phenomenological Consequences}

This section records observationally suggestive consequences of compact $\Sthree$ geometry and
$\golden$-hierarchy phenomenology.  They are not topology detections.

\subsection{Low-multipole suppression}

A compact spatial section has a largest global wavelength scale.  For a round $\Sthree$ of
radius $R_{0}$, scalar harmonics have Laplacian eigenvalues
\begin{equation}
-\Delta Y_{n\ell m}=\frac{n(n+2)}{R_{0}^{2}}Y_{n\ell m},
\qquad n=0,1,2,\ldots,
\label{eq:laplace-spectrum}
\end{equation}
with degeneracy
\begin{equation}
g_{n}=(n+1)^{2}.
\end{equation}
Thus the infrared spectrum is discrete.  This motivates possible suppression or distortion of
large-angle CMB power.

\begin{ansatz}[Low-$\ell$ suppression]
For a compact $\Sthree$ model with radius comparable to the last-scattering diameter, one may
parametrize the low-$\ell$ spectrum by
\begin{equation}
C_{\ell}=C_{\ell}^{\mathrm{ref}}\,
\mathcal{F}_{\ell}(R_{0},\mathcal{T}),
\end{equation}
where $\mathcal{T}$ denotes the transfer calculation and
$\mathcal{F}_{\ell}\to1$ in the large-radius limit.
\end{ansatz}

\begin{remark}
A genuine prediction requires a Boltzmann transfer calculation in the $\Sthree$ harmonic basis.
Low-$\ell$ anomalies alone are not evidence for $\Sthree$ topology.
\end{remark}

\subsection{Log-periodic power-spectrum modulation}

If the $\golden$-hierarchy is imprinted on primordial or late-time perturbations, the simplest
phenomenological modulation is
\begin{equation}
P(k)=P_{\mathrm{sm}}(k)\left[1+A\cos\left(\omega\ln\frac{k}{k_{0}}+\delta\right)\right],
\qquad
\omega=\frac{2\pi}{\ln\golden}.
\label{eq:logperiodic}
\end{equation}

\begin{ansatz}[Golden-ratio log-periodicity]
Equation \eqref{eq:logperiodic} is the minimal log-periodic ansatz associated with the discrete
scale ratio $\golden$.
\end{ansatz}

\begin{remark}
The frequency $\omega=2\pi/\ln\golden$ is sharp if the ansatz is adopted.  The amplitude $A$,
phase $\delta$, pivot $k_{0}$, and physical origin of the modulation are not fixed by topology.
\end{remark}

\subsection{Compact-mode versus box-mode simulations}

A periodic Euclidean box uses $T^{3}$ boundary conditions.  Its eigenmodes are Fourier modes
with an Abelian lattice structure.  A round $\Sthree$ uses hyperspherical harmonics with
spectrum \eqref{eq:laplace-spectrum}.  Therefore a periodic-box simulation does not directly
test the $\Sthree$ mode basis, even when both spaces are compact.

\section{\texorpdfstring{$\Szero$}{S0}, Dependent Distinction, and the Chain of Carriers}
\label{sec:s0-carriers}

This section develops the lowest level of the ladder \eqref{eq:suspension-ladder} under the same
claim-status discipline as the rest of the paper.  The propositions below are elementary
topology; every reading attached to them is explicitly labelled \emph{Interpretation} and proves
no physics.  The starting point of the programme is not an assertion of emergence from nothing:
no absolute arising is assumed anywhere in this section.

\subsection{Minimal distinction and its non-contractibility}

$\Szero=\{-1,+1\}$ should not be pictured as two particles in space.  At level~$0$ no metric,
no phase, no direction, and no causal structure has been introduced.  $\Szero$ is the minimal
two-pole structure of formal distinguishability: this / not-this.  One point alone does not carry
the structure; the two points co-define it.

\begin{proposition}[Non-contractibility of minimal distinction]
\label{prop:s0-noncontract}
\leavevmode
\begin{enumerate}[label=(\roman*)]
\item $\Szero$ is not contractible; indeed $\widetilde H_{0}(\Szero)\simeq\Z$.
\item Removing either point yields a one-point space, which no longer carries the distinction.
\item The identification map $q:\Szero\to\{*\}$ is a quotient, not a homotopy equivalence:
it annihilates $\widetilde H_{0}$.
\item There is no map from the nonempty space $\Szero$ to $\varnothing$.
\end{enumerate}
Consequently, within operations preserving homotopy type, the two-component distinction
persists; it can be destroyed only by an explicit operation of deletion or identification, never
by continuous deformation.
\end{proposition}

\begin{proof}
(i) $\Szero$ has two path components, so $H_{0}(\Szero)\simeq\Z\oplus\Z$ and
$\widetilde H_{0}(\Szero)\simeq\Z$; a contractible space has vanishing reduced homology.
(ii) is immediate.  (iii) The one-point space has $\widetilde H_{0}=0$, and a homotopy
equivalence induces isomorphisms on reduced homology.  (iv) A map must assign a value to every
point of its domain, and $\varnothing$ contains none.
\end{proof}

\begin{remark}[Interpretation: interdependence cannot be contracted, only failed to be preserved]
The two points of $\Szero$ have no meaning \emph{as} $\Szero$ separately: each is a pole of
distinction only because the other exists.  Proposition~\ref{prop:s0-noncontract} is the
mathematical trace of this interdependence.  If a description asserts that the distinction has
``disappeared,'' it must exhibit the operator responsible---a quotient, a deletion, a change of
category---and pay its price: the erasure of the invariant $\widetilde H_{0}$.  Nothing in the
structure-preserving language \emph{derives} that disappearance; it can only be imposed as an
operation of forgetting.
\end{remark}

\subsection{Suspension carries the distinction upward}

\begin{proposition}[The fundamental class of $\Sthree$ is a suspended distinction]
\label{prop:susp-class}
The suspension isomorphism $\widetilde H_{k}(X)\simeq\widetilde H_{k+1}(\Sigma X)$, applied
three times along $\Szero\to\Sone\to\Stwo\to\Sthree$, gives
\begin{equation}
\widetilde H_{0}(\Szero)\;\simeq\;\widetilde H_{1}(\Sone)\;\simeq\;
\widetilde H_{2}(\Stwo)\;\simeq\;\widetilde H_{3}(\Sthree)\;\simeq\;\Z,
\label{eq:susp-transport}
\end{equation}
carrying the distinction class of $\Szero$ to a generator of the fundamental class of
$\Sthree$.
\end{proposition}

\begin{proof}
The suspension isomorphism holds for every space and every degree, by Mayer--Vietoris applied
to the two cones of $\Sigma X$~\cite{Hatcher2002}.  Along $\Sigma\Sn{n}\simeq\Sn{n+1}$ it
identifies generators up to sign, giving \eqref{eq:susp-transport}.
\end{proof}

\begin{remark}[Interpretation: carrier changes; generative relation persists]
Suspension does not create the invariant from non-being and does not destroy it: it re-presents
it on a richer carrier.  What is bare polarity on $\Szero$ acquires phase on $\Sone$, direction
on $\Stwo$, and the spin/Hopf carrier on $\Sthree$.  In the vocabulary of the programme:
\emph{information is not produced from non-being; it is transformed through carriers of
distinguishability}.  A new state of comprehension does not entail a new state of the invariant
content; the invariant is preserved while its mode of representation changes.  Minimal
distinction is not annihilated by suspension; it is carried upward.  Destruction of the
distinction would require an additional quotient operation, not a homotopy.  (This is a
statement about the formal model, not a physical conservation law.)
\end{remark}

\subsection{Loop-trivial, carrier-nontrivial}

\begin{proposition}[Geometric consequence]
\label{prop:loop-carrier}
$\pi_{1}(\Sthree)=0$, while $H_{3}(\Sthree;\Z)\simeq\Z$.  Hence every loop in $\Sthree$
contracts, but $\Sthree$ itself is not contractible; in particular it is not homotopy equivalent
to the ball $B^{3}$, which is contractible and has boundary $\Stwo$, whereas $\Sthree$ is
compact and boundaryless.
\end{proposition}

The short formula is: on $\Sthree$ one can contract every loop, but the carrier itself cannot be
contracted to a point without destroying its fundamental class.  Throughout this paper
$\Sthree$ is therefore used as a compact, boundaryless, simply connected carrier with trivial
fundamental group, \emph{not} as a contractible space:
\begin{quote}
$\Sthree$ is \emph{loop-trivial but carrier-nontrivial}:
$\pi_{1}(\Sthree)=0$, while $H_{3}(\Sthree;\Z)\simeq\Z$.
\end{quote}
At the bottom of the chain the same pattern already holds in degree zero:
$\widetilde H_{0}(\Szero)\simeq\Z$ is the trace of the distinction of components, and it is
precisely this class that \eqref{eq:susp-transport} transports upward.

\subsection{Planck realization of distinguishability}

\begin{remark}[Interpretation: the lower operational grain]
$\Szero$ as a topological object carries no distance.  But if minimal distinguishability
receives a physical-geometric realization inside an effective theory, the natural lower
operational scale of distinction is the Planck grain,
\begin{equation}
\ell_{P},\qquad t_{P}=\ell_{P}/c .
\end{equation}
The exact formula is: \emph{$\Szero$ is the formal carrier of minimal distinction; $\ell_{P}$
and $t_{P}$ are the lower metric-causal realization of distinguishability in effective physical
reconstruction.}  Below this grain the effective model is not obliged to assign an operationally
verifiable causal-geometric state.  This does not mean that reality is absent there; it means
that the language of distinguishable causal-geometric reconstruction has reached its lower
boundary.  Without scale there is no distance; without time there is no causal relation; without
causal relation there is no physically unfolded information; without distinguishability there is
no observable phenomenon.  In brief: distinction $\to$ carrier $\to$ order/phase $\to$ causal
relation $\to$ reconstructible information.
\end{remark}

\begin{remark}[Interpretation: connection with Ricci-type reconstruction]
In the companion Ricci-flow work~\cite{BateninPreece2026} the same boundary appears
dynamically: a reflecting barrier $a\ge\ell_{P}$ defines a forward-invariant region with
$|\mathrm{Rm}|^{2}\le C/\ell_{P}^{4}$, mirroring a self-adjoint Wheeler--DeWitt endpoint on
$a\in[\ell_{P},\infty)$ with zero probability flux through $a=\ell_{P}$.  The flow there is not
the literal time of reality but an instrument of reconstruction: it shows which geometric modes
smooth out, which lead to blow-up, where forcing is insufficient, and where a boundary condition
becomes necessary.  The reading fixed by that programme applies verbatim here:
\begin{quote}
\emph{The Planck boundary is not the beginning of the universe; it is the universal lower
boundary of quasi-local geometric reconstruction.}
\end{quote}
\end{remark}

\subsection{\texorpdfstring{$\Sfour$}{S4} as another representational regime}

$\Sfour$ is not a simple continuation of $\Sthree$ in the same sense.  It remains a compact,
boundaryless, simply connected sphere, but it is not $\SU(2)$, not the group of unit
quaternions, and not parallelizable ($\chi(\Sfour)=2$, so every smooth tangent vector field has
a zero); and it is the \emph{base}, not the total space, of the quaternionic Hopf fibration
$\Sthree\hookrightarrow\Sseven\twoheadrightarrow\Sfour$.  A terminologically precise description
is the one-point compactification of the quaternionic line,
\begin{equation}
\Hbb\cup\{\infty\}\;\simeq\;\HP^{1}\;\simeq\;\Sfour ,
\end{equation}
consistent with Appendix~A.  \emph{Interpretation:} $\Sfour$ marks a transition to a regime in
which the preceding conditioned form of distinguishability is transformed rather than continued
as the same phenomenal configuration.  $\Sthree$ is the place of unfolded spin/Hopf
phenomenology; $\Sfour$ is another carrier-class, where that phenomenology is not simply
prolonged but re-represented.  It is not destroyed, but it ceases to be the same conditioned
configuration.

\subsection{Two protections: against self-being and against substantializing dependence}

\begin{remark}[Interpretation: reading discipline]
Two complementary protections govern the ontological reading of the chain; they are a reading
discipline, not physics.  The first protects against the false view of self-existing arising
(\emph{svabh\=ava}): if self-existing arising is asserted, it must survive the classical
tetralemma of origination---from itself, from another, from both, or without
cause~\cite{Garfield1995}---and it survives none of them.  The second protects dependence
itself from being turned into a new substance: after refuting self-being one may not posit a
special self-existing entity called ``Interdependence,'' which would be the same error under
another name.  The compressed formula is: \emph{the first part cuts off the self-being of the
phenomenon; the second part cuts off the self-being of dependence itself.}  In mathematical
register: the object is not self-sufficient; it is given by a construction.  It is empty of
own-nature, but not empty of dependent functioning.

The same discipline applies to space, matter, and information.  Space is neither denied nor
absolutized; it is considered as a dependently designated structure disclosed through matter,
causal relations, and observer reconstruction.  Information is neither a self-existing thing nor
nothing; it is dependently designated through distinguishability, carrier, scale, causality, and
the possibility of reconstruction.
\end{remark}

\subsection{Suspension hierarchy principle}

\begin{tcolorbox}[colback=gray!4,colframe=gray!50,title=Suspension hierarchy principle (Interpretation)]
A minimal distinction represented by $\Szero$ may be transformed through a sequence of carriers
$\Sigma\Sn{n}\simeq\Sn{n+1}$.  At each step the invariant content is neither simply identical
nor simply different: it is preserved through a change of representational carrier
(Proposition~\ref{prop:susp-class}).  The endpoint is not annihilation of information, but a
regime in which the conditioned mode of distinction has been transformed.
\end{tcolorbox}

In brief: $\Szero$ gives minimal distinguishability; suspension changes the carriers of this
distinguishability; $\Sthree$ gives the special compact spin/Hopf carrier; the Planck scale
gives the lower metric-causal boundary of physical reconstruction; interdependence explains
functioning but does not itself become a substance.  Not self-being, not nothing, but
dependently reconstructible distinguishability passing through a chain of carriers.

\section{Conclusion}

The suspension ladder
\[
\Szero\to\Sone\to\Stwo\to\Sthree\to\Sfour
\]
has a precise mathematical reading.  $\Szero$ supplies binary distinction; $\Sone$ supplies
phase; $\Stwo$ supplies direction; $\Sthree$ is the first sphere where phase and direction become
a nontrivial total space carrying $\SU(2)$, spin geometry, parallelizable flow, and compact
spectral discreteness; $\Sfour\simeq\HP^{1}$ supplies a quaternionic compensation endpoint.
At the base of the ladder, $\Szero$ supplies minimal interdependent distinction, transported by
the suspension isomorphism onto the fundamental class of $\Sthree$; the Planck grain is its
lower operational realization in effective reconstruction (Sec.~\ref{sec:s0-carriers}).

The central theorem is not that all physics is solved by topology.  The theorem is the minimal
carrier statement: $\Sthree$ is the first sphere in the chain satisfying the joint spin--Hopf--flow
conditions.  The golden-ratio hierarchy gives a minimal additive self-similar scaling; its entropy
monotonicity follows under an area-counting rule; its recurrence-breaking interpretation remains
conjectural until a Hamiltonian measure is specified.  Low-multipole suppression and
log-periodic modulation are observational ansätze, not detections.

The mathematical core is therefore strong but bounded:
\[
\text{topology supplies the carrier,}
\qquad
\text{dynamics and amplitudes require additional laws.}
\]
This is the clean form of the suspension-of-spheres programme.

\appendix

\section{\texorpdfstring{$\HP^{1}\simeq\Sfour$}{HP1 = S4}}

The quaternionic projective line is
\begin{equation}
\HP^{1}=(\Hbb^{2}\setminus\{0\})/\Hbb^{\times}.
\end{equation}
Every equivalence class has a representative $(q_{1},q_{2})$ with
$|q_{1}|^{2}+|q_{2}|^{2}=1$, unique up to right multiplication by a unit quaternion.
Thus
\[
\Sthree\hookrightarrow\Sseven\twoheadrightarrow\HP^{1}.
\]
An explicit map $\Sseven\to\Sfour\subset\R^{5}$ is
\begin{equation}
(q_{1},q_{2})\longmapsto
\left(
|q_{1}|^{2}-|q_{2}|^{2},\,
2q_{1}\overline{q}_{2}
\right)\in\R\oplus\Hbb\simeq\R^{5}.
\end{equation}
The norm of the image is
\[
(|q_{1}|^{2}-|q_{2}|^{2})^{2}+4|q_{1}|^{2}|q_{2}|^{2}
=(|q_{1}|^{2}+|q_{2}|^{2})^{2}=1,
\]
so the image lies on $\Sfour$.  This gives the standard identification
$\HP^{1}\simeq\Sfour$.

\section{RG Interpretation of the \texorpdfstring{$\golden$}{phi}-Hierarchy}

A cautious renormalization-group reading of the hierarchy assigns a scale
\[
\Lambda_{n}=\Lambda_{0}\golden^{-n}
\]
or, equivalently, a length
\[
R_{n}=R_{0}\golden^{n}.
\]
The hierarchy is then a discrete RG skeleton with step size $\ln\golden$.  A log-periodic
observable has the generic form
\[
F(x)=F_{0}(x)\left[1+A\cos\left(\frac{2\pi\ln x}{\ln\golden}+\delta\right)\right].
\]
This appendix does not claim that the RG flow is derived; it records the minimal form that any
$\golden$-scale phenomenology must take.

\section{Hausdorff Dimension: Heuristic Only}

If an iterated-function system with $N$ self-similar copies at contraction ratio $r$ is specified,
its similarity dimension is determined by
\[
Nr^{d}=1.
\]
Without an explicit IFS, the phrase ``fractal $\Sthree$'' has no unique Hausdorff dimension.
For example, taking $r=\golden^{-1}$ and $N\simeq\golden^{3}$ gives $d\simeq3$, but this is a
heuristic dimensional consistency check, not a theorem.

\section{Additional Observational Notes}

\paragraph{CMB topology searches.}
Circles-in-the-sky and matched-correlation searches constrain compact topologies, but a null
result only excludes compact radii below the search sensitivity and specific quotient structures.

\paragraph{Matter power spectrum.}
The log-periodic form \eqref{eq:logperiodic} gives a sharp frequency target
$\omega=2\pi/\ln\golden$, but absence of this feature would falsify the imprint ansatz, not the
pure topological minimality theorem.

\paragraph{Gravitational waves.}
Tensor modes on compact $\Sthree$ also have discrete harmonic labels.  Turning this into a
detectable gravitational-wave line prediction requires a source model, transfer function, and
instrumental likelihood.  Compactness alone gives the mode basis, not the stochastic amplitude.

\begin{thebibliography}{99}

\bibitem{Hatcher2002}
A.~Hatcher, \emph{Algebraic Topology} (Cambridge University Press, 2002).

\bibitem{Hopf1931}
H.~Hopf, \emph{\"Uber die Abbildungen der dreidimensionalen Sph\"are auf die Kugelfl\"ache},
Math.\ Ann.\ \textbf{104}, 637--665 (1931).

\bibitem{Milnor1997}
J.~Milnor, \emph{Topology from the Differentiable Viewpoint} (Princeton University Press, 1997).

\bibitem{Milnor1956}
J.~Milnor, \emph{On manifolds homeomorphic to the 7-sphere},
Ann.\ of Math.\ \textbf{64}, 399--405 (1956).

\bibitem{Cannon1979}
J.~W.~Cannon, \emph{The recognition problem: what is a topological manifold?},
Bull.\ Amer.\ Math.\ Soc.\ \textbf{84}, 832--866 (1978).

\bibitem{Edwards1980}
R.~D.~Edwards, \emph{The topology of manifolds and cell-like maps},
Proceedings of the International Congress of Mathematicians, Helsinki (1978), 111--127.

\bibitem{LawsonMichelsohn1989}
H.~B.~Lawson and M.-L.~Michelsohn, \emph{Spin Geometry} (Princeton University Press, 1989).

\bibitem{Besse1987}
A.~L.~Besse, \emph{Einstein Manifolds} (Springer, 1987).

\bibitem{Baez2002}
J.~Baez, \emph{The Octonions}, Bull.\ Am.\ Math.\ Soc.\ \textbf{39}, 145--205 (2002).

\bibitem{Perelman2002}
G.~Perelman, \emph{The entropy formula for the Ricci flow and its geometric applications},
arXiv:math/0211159 (2002).

\bibitem{Perelman2003a}
G.~Perelman, \emph{Ricci flow with surgery on three-manifolds},
arXiv:math/0303109 (2003).

\bibitem{Perelman2003b}
G.~Perelman, \emph{Finite extinction time for the solutions to the Ricci flow on certain
three-manifolds}, arXiv:math/0307245 (2003).

\bibitem{Bousso1999}
R.~Bousso, \emph{A covariant entropy conjecture}, JHEP \textbf{07}, 004 (1999).

\bibitem{Planck2018}
Planck Collaboration, \emph{Planck 2018 results. VI. Cosmological parameters},
Astron.\ Astrophys.\ \textbf{641}, A6 (2020).

\bibitem{Cornish1998}
N.~J.~Cornish, D.~N.~Spergel, and G.~D.~Starkman,
\emph{Circles in the sky: finding topology with the microwave background radiation},
Class.\ Quantum Grav.\ \textbf{15}, 2657--2670 (1998).

\bibitem{CMBS4}
CMB-S4 Collaboration, \emph{CMB-S4 Science Book, First Edition}, arXiv:1610.02743 (2016).

\bibitem{DESI2024}
DESI Collaboration, \emph{DESI 2024 VI: cosmological constraints from the measurements of baryon
acoustic oscillations}, arXiv:2404.03002 (2024).

\bibitem{BateninPreece2026}
B.~Batenin and A.~Preece, \emph{Hopf-Lifted Ricci-Type Flow on $\Sthree$: A No-Go for Bounded
and Engineered Forcing, and a Planck-Barrier Invariant Region}, research manuscript (2026).

\bibitem{Garfield1995}
N\=ag\=arjuna, \emph{M\=ulamadhyamakak\=arik\=a}; English translation and commentary:
J.~L.~Garfield, \emph{The Fundamental Wisdom of the Middle Way} (Oxford University Press, 1995).

\end{thebibliography}

\end{document}
